\note{sa page 12}

\struck{\ill{}} $\Phi_{0}$ est un module libre sur $\mathbb{Z}$, engendré par les $\varphi^{\beta,\alpha}_{i}$
\struck{\ill{}} \struck{$L^{\beta} \pi_{i,\alpha}$}
\struck{$\varphi^{\beta}_{i,\alpha}$ \ill{}}, avec

\[ 0 \leqslant i \leqslant n , \quad 0 \leqslant \alpha, \beta \leqslant 2n-i \]

\note{À droite, biffé d'un trait, un premier jeu d'indices : « $0 \leqslant i
\leqslant 2n$, $0 \leqslant j \leqslant 2n$, $\alpha, \beta \geqslant 0$,
$i - 2\alpha = j - 2\beta \geqslant 0$ ».}

\[ \begin{cases} \varphi^{\beta,\alpha}_{i} = L_{0}^{\beta-\alpha}\, \pi_{i+2\alpha,\alpha} & \text{si } \alpha \leqslant \beta \\[2pt] \varphi^{\beta,\alpha}_{i} = \Lambda_{0}^{\alpha-\beta}\, \pi_{i+2\alpha,\alpha} & \text{si } \alpha \geqslant \beta \end{cases} \]
\[ \begin{cases} = L_{0}^{\beta-\alpha}\, \pi_{i,\alpha} = \pi_{j,\beta}\, L_{0}^{\beta-\alpha} & \text{si } \alpha \leqslant \beta \\[2pt] = \Lambda_{0}^{\alpha-\beta}\, \pi_{i,\alpha} = \pi_{j,\beta}\, \Lambda_{0}^{\alpha-\beta} & \text{si } \alpha \geqslant \beta \end{cases} \]

\marginal{Les $\varphi^{\beta,\alpha}_{i}$ et leur multiplication \uncertain{devraient être}
donnés au n° 6}
La multiplication des \uncertain{éléments de cette base} ne dépend pas du choix de
$(M_{0}, \xi_{0})$ (ni du fait \ill{} $\mathbb{Z}$ \ldots), donc l'\ill{}
\ill{} \uncertain{en termes de l'entier} $2n$. $\Phi_{0}$ est bien
déterminé \ill{}.

\textbf{Proposition.} Revenant à la donnée de $M^{*}$ comme au n° 1,
$\mathcal{E}$ \struck{$n - 2g$} \add{anneau} gradué des \struck{\ill{}}
$\operatorname{End}(M^{*})$, et de $L \in \mathcal{E}^{(2)}$.
Conditions équivalentes :

\begin{enumerate}
  \item L'\uncertain{endom.}\ $L$ de $M^{*}$ satisfait la condition de \struck{Hodge}
        Lefschetz, et l'opérateur $\Lambda$ associé est dans $\mathcal{E}$ (i.e.\ on
        est sous les conditions du th.\ du n° 6).
  \item il existe un hom.\ d'anneaux gradués
        $\varphi \colon \Phi_{0} \to \mathcal{E}$
\end{enumerate}

